\subsection{希尔伯特空间的情形}

在本节中，E 表示 C 上的一个希尔伯特空间。我们用 \(\pi\) 表示从 \(\mathcal{L}(\mathrm{E})\) 到卡尔金代数 \(\mathcal{Calk}(\mathrm{E})\) 上的典范同态。

命题 7. --- 设 \(u \in \mathcal{L}(\mathrm{E})\) 。

\begin{enumerate}
\def\labelenumi{\alph{enumi})}
\item
  若 u 是正规的，则 u 是里斯自同态当且仅当 u 是指标为 0 的弗雷德霍姆自同态。
\item
  若 u 是埃尔米特的，则 u 是弗雷德霍姆自同态当且仅当 u 是里斯自同态。
\end{enumerate}

每个里斯自同态都是指标为 0 的弗雷德霍姆自同态（III, p. 46 的命题 5）。反之，假设 u 是指标为 0 的弗雷德霍姆自同态且 u 是正规的。此时它的幂零空间与它的核重合（EVT, V, p. 43, 命题 8），故有限维。于是由 III, p. 45 的引理 2 以及定义，u 是里斯自同态。

我们来证明 \(b\) 。由 \(a\) ，只需验证埃尔米特弗雷德霍姆自同态 \(u\) 的指标为零。但此时 \(u\) 的像（它是闭的）的正交补等于 \(u\) 的核（EVT, V, p. 41, 命题 4），由此得该断言。

推论. --- 设 \(u \in \mathcal{L}(\mathrm{E})\) 。若 u 是正规的，则 \(\mathrm{Sp}_{0}(u)\) 是空的；若 u 是埃尔米特的，则 \(\mathrm{Sp}_{\mathrm{e}}(u)\) 与 \(\pi(u)\) 相对于代数 \(\mathcal{Calk}(\mathrm{E})\) 的谱重合。

两个断言都由该命题以及诸集 \(\mathrm{Sp}_{n}(u)\) （\(n \in \overline{Z}\) ）与 \(\mathrm{Sp}_{\omega}(u)\) 的定义得出，它们构成 u 的本质谱的一个分划（III, p. 83 的定义 2）。

定理 3（外尔）. --- 设 \(u \in \mathcal{L}(\mathrm{E})\) 是 E 的一个正规自同态。\(u\) 的本质谱是诸集 \(\mathrm{Sp}(u + h)\) （当 \(h\) 取遍 \(\mathcal{L}^{\mathrm{c}}(\mathrm{E})\) 时）的交。

设 \(h \in \mathcal{L}^c(\mathrm{E})\) 。由于 \(\mathrm{Sp}_0(u)\) 是空的（上面的推论），III, p. 89 的定理 2 蕴含 \(\mathrm{Sp_e}(u + h) = \mathrm{Sp_e}(u)\) 。于是诸集 \(\mathrm{Sp}(u + h)\) 的交包含 \(\mathrm{Sp_e}(u)\) 。

设 \(\lambda \in \mathrm{Sp}_{\mathrm{s}}(u)\) 。用 \(\mathrm{E}_{\lambda}\) 表示 \(u\) 的相对于 \(\lambda\) 的特征空间，用 \(\mathrm{F}_{\lambda}\) 表示与 \(u\) 及 \(\mathbf{C} = \{\lambda\}\) 相伴的谱投影的像。空间 E 是 \(\mathrm{E}_{\lambda}\) 与 \(\mathrm{F}_{\lambda}\) 的拓扑直和。设 \(h\) 是 E 的这样一个有限秩自同态：它在 \(\mathrm{F}_{\lambda}\) 上为零，而在 \(\mathrm{E}_{\lambda}\) 上与恒同映射重合。自同态 \(u + h\) 可逆，故 \(\lambda \notin \mathrm{Sp}(u + h)\) 。定理由此得出。

于是，若 u 是 E 的正规自同态，且 \(h \in \mathcal{L}^{c}(\mathrm{E})\) ，则 \(u + h\) 的谱与 u 的谱只能相差一些具有有限谱重数的孤立点。
% semantic-fragments: legacy-input-seam
\relax{}
